\paragraph{4.2. Алгоритм Штрассена для умножения матриц}

\subparagraph{4.2.1}
Воспользуйтесь алгоритмом Штрассена для вычисления произведения матриц
\[
\begin{pmatrix}
1 & 3 \\
7 & 5
\end{pmatrix}
\begin{pmatrix}
6 & 8 \\
4 & 2
\end{pmatrix}.
\]
Покажите, как вы это делаете.

\textbf{Ответ:}


\subparagraph{4.2.2}
Запишите псевдокод алгоритма Штрассена.

\textbf{Ответ:}
\begin{lstlisting}
class StrassenMatrixMultiplyRecursive
{
public:
    template <typename T, size_t N>
    Matrix<N,T> operator()(const Matrix<N,T>& A, const Matrix<N,T>& B)
    {
        Matrix<N,T> C;
        multiply(A, B, S, C);
        return C;
    }

private:
    template <typename T, size_t N>
    void multiply(const DynamicMatrixWrapper& A, const DynamicMatrixWrapper& B, DynamicMatrixWrapper& C)
    {
        auto aSize = A.size();
        if (aSize == 1)
        {
            C[0][0] = A[0][0] * B[0][0];
            return;
        }

        auto A11 = A.SubMatrix({0,0},{aSize/2, aSize/2});
        auto A12 = A.SubMatrix({0,aSize/2},{aSize/2, aSize});
        auto A21 = A.SubMatrix({aSize/2,0},{aSize, aSize/2});
        auto A22 = A.SubMatrix({aSize/2, aSize/2},{aSize, aSize});

        auto bSize = B.size();
        auto B11 = B.SubMatrix({0,0},{aSize/2, aSize/2});
        auto B12 = B.SubMatrix({0,aSize/2},{aSize/2, aSize});
        auto B21 = B.SubMatrix({aSize/2,0},{aSize, aSize/2});
        auto B22 = B.SubMatrix({aSize/2, aSize/2},{aSize, aSize});
        
        std::array<DynamicMatrix, 10>& S {aSize, bSize};
        S[0] = B12 - B22;
        S[1] = A11 + A12;
        S[2] = A21 + A22;
        S[3] = B21 - B11;
        S[4] = A11 + A22;
        S[5] = B11 + B22;
        S[6] = A12 - A22;
        S[7] = B21 + B22;
        S[8] = A11 - A21;
        S[9] = B11 + B12; 
        ....
    }
}
\end{lstlisting}

\subparagraph{4.2.3}
Как модифицировать алгоритм Штрассена для перемножения матриц размером \(n \times n\),
где \(n\) не является точной степенью 2? Покажите, что получающийся в результате алгоритма
выполняется за время \(O(n^{\lg 7})\).

\textbf{Ответ:}

Дополняем матрицы нулями до ближайшей степени 2, большей или равной \(n\), 
затем применяем стандартный алгоритм Штрассена. 
Пусть \(N = 2^{\lceil \log_2 n \rceil} < 2n\). 
Время работы \(T(N) = \Theta(N^{\lg 7}) = O((2n)^{\lg 7}) = O(n^{\lg 7})\).


\subparagraph{4.2.4}
Чему равно наибольшее \(k\), такое, что если вы можете перемножить \(3 \times 3\)-матрицы с
помощью \(k\) умножений (не предполагая коммутативности умножения), 
то вы можете перемножить матрицы размером \(n \times n\) за время \(O(n^{\lg 7})\)? 
Каким должно быть время работы такого алгоритма?

\textbf{Ответ:}



\subparagraph{4.2.5}
В. Пан (V. Pan) открыл способ перемножения матриц размером \(68 \times 68\) с использованием только 132.464 умножений, 
способ перемножения матриц размером \(70 \times 70\) с использованием 143.640 умножений 
и способ перемножения матриц размером \(72 \times 72\) с использованием 155.424 умножений. 
Какой из методов дает нам лучшее асимптотическое время работы при его использовании в алгоритме 
"разделяй и властвуй" для перемножения матриц? Проведите сравнение с алгоритмом Штрассена.

\subparagraph{4.2.6}
Насколько быстро вы сумеете умножить матрицу размером \(kn \times n\) на матрицу размером \(n \times kn\), применяя алгоритм Штрассена в качестве подпрограммы? Ответьте на тот же вопрос для ситуации, когда мы меняем входные матрицы местами.

\subparagraph{4.2.7}
Покажите, как перемножить комплексные числа \(a + bi\) и \(c + di\), 
используя только три умножения действительных чисел. 
Алгоритм должен получать \(a\), \(b\), \(c\) и \(d\) 
в качестве входных данных и возвращать действительную (\(ac - bd\)) 
и мнимую (\(ad + bc\)) части произведения по отдельности.